%\section{Lexical Semantic Change and Concept Creep}

Lexical semantic change (LSC) refers to shifts in a word's meaning while its grammatical function remains stable, and constitutes a common form of language change \citep{campbellHistoricalLinguisticsIntroduction2013}. For example, \textit{mouse}, initially a word for a small rodent, broadened in usage to refer to a computer input device.

\textit{Concept creep theory} posits that harm-related concepts are particularly prone to LSC, specifically gradual semantic broadening. As such, many harm-related terms, including \textit{addiction}, \textit{bullying}, \textit{harassment}, \textit{prejudice}, and \textit{trauma}, are reported to have expanded in meaning since at least the 1970s \citep{haslamConceptCreepPsychologys2016,baesSemanticInflationTrauma2023}. Concept creep distinguishes between two forms of LSC that can occur concurrently: vertical creep and horizontal creep. Vertical creep refers to the loosening of definitions to include milder instances (declining semantic severity or intensity), whereas horizontal creep refers to the extension of definitions to encompass qualitatively new phenomena (semantic broadening). Concepts of mental illness are considered harm-related because distress and dysfunction are fundamental to their definition \citep{wakefieldConceptMentalDisorder1992}. Previous studies have, therefore, examined concept creep across different mental-health-related concepts.

\citet{vylomovaSemanticChangesHarmrelated2021} found that \textit{trauma} and \textit{addiction} (together with three non-mental health-related concepts: \textit{bullying}, \textit{prejudice}, and \textit{harassment}) exhibited both vertical and horizontal creep in a corpus of approximately 800,000 psychology article abstracts from 875 journals dating back to the 1960s. Similar, though weaker, patterns were observed in a general-language corpus combining the Corpus of Contemporary American English (CoCA) and the Corpus of Historical American English (CoHA) from the 1970s to the 2010s. Subsequent studies using the same corpora and time periods have reported comparable findings for other mental-health-related concepts. \citet{baesSemanticInflationTrauma2023} found evidence of vertical creep in \textit{trauma} in the psychology corpus. \citet{baesSemanticShiftsMental2023} reported vertical creep in \textit{addiction}, \textit{anger}, \textit{stress}, and \textit{worry} in the psychology corpus, and in \textit{addiction}, \textit{grief}, \textit{stress}, and \textit{worry} in the general corpus. \citet{baesMultidimensionalFrameworkEvaluating2024} further found that the broader concepts \textit{mental health} and \textit{mental illness} both expanded horizontally in the psychology corpus.

Not all findings align with this pattern. One study reported that the semantic intensity of \textit{anxiety} and \textit{depression} increased in both Vylomova and Haslam's psychology corpus and their general-language corpus, contrary to concept-creep expectations \citep{xiaoHaveConceptsAnxiety2023}. This unexpected result may reflect measurement unspecificity, particularly the absence of discourse-context and construct distinctions. For example, \textit{depression} may refer to psychiatric depression, but also to meteorological or economic phenomena. Similarly, \textit{anxiety} may refer to a nosological category, a disorder construct, or an underlying human experience. In a replication study, \citet{pislRevisitingSemanticSeverity2025} showed that this apparent increase in semantic intensity could be attributed to changing discourse composition, specifically shifts in the balance between clinical or nosological contexts and lived-experience contexts, rather than to intrinsic semantic change alone. In their analysis of lead paragraphs from \textit{New York Times} articles published between 1970 and 2023, the time effect for \textit{depression} became nonsignificant after controlling for mental-health context.

Another study examined a range of terms commonly associated with therapy-speak, including \textit{toxic}, \textit{bipolar}, \textit{psychopath}, \textit{narcissistic}, and \textit{triggered}, and reported mixed results \citep{iacobComputationalAnalysisEmergence2026}. In an extension of Vylomova and Haslam's psychology corpus, most terms exhibited horizontal creep. By contrast, in two Reddit corpora comprising comments from psychology-oriented subreddits and general subreddits between 2010 and 2025, most terms showed a narrowing in breadth. Overall, the authors found that long-established psychological terms such as \textit{OCD}, \textit{bipolar}, and \textit{trauma} displayed little semantic change, whereas terms such as \textit{gaslighting} and \textit{imposter} shifted substantially from year to year.

On balance, many mental-health-related concepts appear to have undergone concept creep, becoming milder and broader over recent decades. Haslam and colleagues theorise several drivers of this process \citep{haslamHarmInflationMaking2020,haslamConceptCreepMental2026}. One is the influence of ``opprobrium entrepreneurs'', who seek to cast previously accepted conditions in a more problematic light; by broadening a harm concept, the disapproval associated with its original meaning can come to apply to less severe cases. A second driver is ``prevalence-induced conceptual change'', whereby the standards for recognising harm tend to loosen as instances of harm become less common. Finally, Haslam and colleagues point to a broader cultural increase in concern with harm. As harm concepts expand, a wider range of experiences is treated as morally significant and worthy of care, and a greater number of actions come to be seen as harmful \citep{haslamHarmInflationMaking2020,haslamConceptCreepMental2026}.

Whether the broadening of mental-health constructs reflects changes in official diagnostic criteria is unclear. A meta-analysis showed that no revision of the \textit{Diagnostic and Statistical Manual of Mental Disorders} (DSM) from the third edition onward was reliably more inflationary or deflationary overall. However, specific disorders changed significantly. Most notably, ADHD inflated by 18\%, 33\%, and then 17\% across the three revisions following DSM-III. Autism also inflated by 50\% from DSM-III to DSM-III-R, albeit based on a single study, but deflated by 15\% from DSM-IV to DSM-5 \citep{fabianoDiagnosticInflationDSM2020}.

\section{Computational Approaches to Studying \texorpdfstring{\\}{ } Lexical Semantic Change}

Semantic change has long been studied across linguistics and the social sciences, but lexical semantic change remains difficult to characterise because changes in word meaning are often gradual and less visible than other forms of linguistic change, such as those produced by spelling or grammar reforms \citep{saSurveyCharacterizingSemantic2026}. Earlier research largely depended on manual methods, with linguists using historical texts, dictionaries, and corpora to reconstruct how word meanings changed over time. More recently, computational linguistics and natural language processing (NLP) have made it possible to study semantic change at much larger scales \citep{jm3}. Within NLP, meaning is typically understood in distributional terms: a word's meaning is inferred from the contexts in which it appears. These contextual patterns can be represented through several computational approaches, including frequency-based measures, topic models, semantic graphs, and embedding-based methods \citep{saSurveyCharacterizingSemantic2026}.

Although considerable progress has been made in identifying LSC using these techniques \citep[see][]{tahmasebiSurveyComputationalApproaches2019,kutuzovDiachronicWordEmbeddings2018,tangStateoftheartSemanticChange2018}, less attention has been paid to characterising the nature of such changes. To address this gap, \citet{baesMultidimensionalFrameworkEvaluating2024} proposed the SIBling framework---Sentiment, Intensity, and Breadth---as a unified, multidimensional approach to characterising LSC. The framework situates concept creep within a broader account of LSC and links this conceptual model to computational methods, including frequency-based and embedding-based techniques. It distinguishes three dimensions along which terms may change over time: vertical drift, horizontal drift, and sentiment. Vertical drift refers to changes in intensity. A term may strengthen, as in \textit{hilarious} shifting from cheerful or amusing to extremely funny, or weaken, as in \textit{trauma} shifting from brain injuries to milder events such as business loss. This dimension maps onto vertical concept creep. Horizontal drift refers to changes in breadth. A term may narrow, as in \textit{doctor} shifting from scholar or teacher to primarily denoting a medical professional, or broaden, as in \textit{cloud} shifting from a meteorological term to internet-based data storage. This dimension maps onto horizontal concept creep. Sentiment refers to changes in connotation. A term may acquire a more positive connotation, as in \textit{geek} shifting from a derogatory term for odd people to someone passionate about a field, or a more negative connotation, as in \textit{retarded} shifting from a neutral term for intellectual disability to a highly pejorative slur. This dimension roughly corresponds to destigmatisation and stigmatisation, which are not directly captured by concept creep theory. According to SIBling, these dimensions can be complemented by examining shifts in target-word frequency, or salience, and in the thematic content of target-word collocates. Despite its conceptual value, however, the framework's operationalisations should be treated as first implementations rather than settled measures: they have not yet been extensively validated externally, and their measurement choices remain open to refinement \citep{baesMultidimensionalFrameworkEvaluating2024}.

\section{The Present Study}

Three gaps in the existing literature motivate the present study.

The first concerns the evidentiary basis of prior work. Research on LSC in mental-health-related concepts has relied heavily on a small set of curated corpora, particularly COCA/COHA and Vylomova and Haslam's \citeyear{vylomovaSemanticChangesHarmrelated2021} psychology abstracts datasets \citep[e.g.,][]{baesSemanticShiftsMental2023,baesMultidimensionalFrameworkEvaluating2024,xiaoHaveConceptsAnxiety2023,iacobComputationalAnalysisEmergence2026}. This matters because the broader aim of this literature is to infer cultural dynamics from lexical change. Yet curated corpora may partly reflect shifts in editorial policy, audience composition, disciplinary conventions, or ideological stance rather than broader changes in public discourse \citep{pislRevisitingSemanticSeverity2025}. To date, LSC in mental-health-related concepts has not been examined systematically in general web discourse.

The second concerns the target concepts themselves. Direct evidence on LSC in ADHD and autism remains limited. One study found that the two concepts have converged semantically on Reddit: from 2019 onward, their contextual similarity increased, and the terms became more similar to each other than to comparison conditions \citep{kangConvergingRepresentationsAttentionDeficit2025}. This finding, however, captures convergence between ADHD and autism on a single platform; it does not characterise how either concept's semantic profile has evolved over time in general web discourse.

The third concerns discourse composition. Existing work indicates that apparent trends in LSC may partly reflect shifts in the contexts in which terms are used—such as clinical categories, service labels, identity or lived-experience categories, and objects of public debate—rather than intrinsic semantic change \citep{pislRevisitingSemanticSeverity2025,xiaoHaveConceptsAnxiety2023}. Treating all mentions as a single semantic population therefore risks conflating genuine semantic change with change in discourse composition.

Taken together, these gaps point to three corresponding requirements: a broader evidentiary base that captures general web discourse, direct diachronic analysis of ADHD and autism, and an analytic approach that accounts for variation in discourse framing. The present study was designed to meet these requirements.

%\subsection{Research Questions and Hypotheses}

Guided by these requirements, the present study addresses four research questions:

\begin{description}
	\item[RQ1:] How has the salience of ADHD and autism in general web discourse changed from 2014 to 2026, against the backdrop of non-clinical negative-emotion baseline terms?
	\item[RQ2:] How has the balance between clinical and lived-experience framing of ADHD and autism changed over this period?
	\item[RQ3:] How have the intensity, breadth, and sentiment of ADHD and autism contexts changed—overall and within clinical versus lived-experience frames—against the backdrop of the baseline terms?
	\item[RQ4:] How has the thematic content of ADHD and autism discourse evolved over time?
\end{description}

Consistent with concept creep theory, it is hypothesised that ADHD and autism will increasingly appear in less emotionally intense contexts (vertical concept creep) and across a broader qualitative range of lexical contexts (horizontal concept creep). This hypothesis corresponds to the intensity and breadth components of RQ3; the remaining measures—sentiment (RQ3), frame composition (RQ2), and thematic content (RQ4)—are treated as exploratory, given the limited prior research available to motivate directional predictions.

%\subsection{Methodological Overview}

Methodologically, the present study adapts Baes et al.'s \citeyear{baesMultidimensionalFrameworkEvaluating2024} SIBling framework for LSC analysis, refining its operationalisation through target-aware embeddings for breadth and an expanded affective lexicon for intensity and sentiment. To broaden the evidentiary base and support direct diachronic semantic analysis of ADHD and autism, the study constructs a diachronic corpus of general web discourse using Common Crawl, a large-scale, openly accessible repository of web crawl data. The corpus spans 2014--2026, from the earliest period with sufficiently consistent and usable data to the most recent available year. An efficient and reproducible pipeline extracts quality-filtered web documents containing ADHD, autism, and baseline emotion terms, supporting both frequency estimation and downstream semantic analysis while ensuring comparability across targets and baselines. Finally, to account for variation in discourse framing, the study incorporates a supervised frame-classification step that distinguishes clinical or disorder-construct discourse from lived-experience discourse prior to LSC analysis.

%\subsection{Contributions}

In sum, the present study contributes to the literature by extending concept creep and therapy-speak research to ADHD and autism; shifting the evidentiary base for LSC research from curated corpora toward general web discourse; introducing a frame-aware approach for separating clinical and lived-experience uses of mental health terms; implementing a reproducible Common Crawl pipeline for LSC research; and assessing whether selected operational refinements can strengthen the application of the SIBling framework.
